\subsection{酉表示的类}

引理 9. --- 设 E 是 K 上的希尔伯特空间，F ⊂ E 是 E 的稠密向量子空间。那么 E 的希尔伯特维数小于或等于 F 的维数。

若 E 的希尔伯特维数有限，则它等于 E 的维数。设 E 的希尔伯特维数无限，则子空间 F 的维数也是无限的。设 B 是 E 的标准正交基，B' 是 F 的基。对于每个 \(x \in B\)，存在 \(f(x) \in B'\)，使得 \(\langle x \mid f(x) \rangle \neq 0\)，否则将有 \(x \in F^{\circ} = \{0\}\)。于是定义映射 \(f: B \to B'\)。

对于每个 \(y \in B'\)，集合 \(f^{-1}(y)\) 包含于满足 \(x \in B\) 且 \(\langle x | y \rangle \neq 0\) 的集合中，因而是可数的（EVT, V, p. 21, 命题 4）。由 E, III, p. 50, 命题 4，得到 \(\text{Card(B)} = \text{Card}(f(B)) \leqslant \text{Card(B')}\)。

记 \(Is_{G}(\pi_{1},\pi_{2})\) 为关系

“G 是拓扑群，且 \(\pi_{1}\) 和 \(\pi_{2}\) 是

K 上希尔伯特空间中的酉表示，并且二者

同构。”关于 \(\pi_{1}\) 和 \(\pi_{2}\)，这是一个等价关系。对于 K 上希尔伯特空间中 G 的任意酉表示 \(\pi\)，记 \(\mathrm{cl}(\pi)\) 为 \(\pi\) 的等价类（参见 E, II, p. 47）；因此它是 G 的一个与 \(\pi\) 同构的酉表示；称 \(\mathrm{cl}(\pi)\) 为 \(\pi\) 的类。K 上希尔伯特空间中的酉表示 \(\pi_{1}\) 和 \(\pi_{2}\) 同构，当且仅当 \(\mathrm{cl}(\pi_{1}) = \mathrm{cl}(\pi_{2})\)。

设 G 是拓扑群，c 是一个基数。关系

“λ 是 G 在

希尔伯特维数 \(\leqslant\mathfrak{c}\) 的复希尔伯特空间中的酉表示的类”

是关于 \(\lambda\) 的集化关系（E, II, p. 3）。事实上，每个维数 \(\leqslant \mathfrak{c}\) 的 K 上希尔伯特空间都等距同构于 \(\ell^2 (\mathfrak{c})\) 的一个希尔伯特子空间（EVT, V, p. 23, cor. 2），于是断言由 E, II, p. 47 得出。

设 \(\pi\) 是 G 在希尔伯特空间 E 中的不可约酉表示。取 E 中的非零元素 \(x\)。由于 \(\pi\) 不可约，向量 \(x\) 是 \(\pi\) 的循环向量，这意味着 E 的希尔伯特维数 \(\leqslant\) Card(G)（将引理 9 应用于由 \(\pi(g)x\)（\(g \in G\)）所生成的稠密子空间）。因此，G 在 K 上希尔伯特空间中的不可约酉表示的各类，属于 G 在 K 上希尔伯特维数 \(\leqslant\) Card(G) 的希尔伯特空间中的酉表示的类所成的集合；所以它们构成一个集合。

定义 8. --- 我们用 \(\widehat{\mathbf{G}}\) 表示 \(\mathbf{G}\) 在复希尔伯特空间中的不可约酉表示的类所成的集合。称 \(\widehat{\mathbf{G}}\) 为 \(\mathbf{G}\) 的酉对偶。

因此，对于 G 在复希尔伯特空间中的每个不可约酉表示 \(\pi\)，都有 \(\operatorname{cl}(\pi) \in \widehat{\mathbf{G}}\)。

备注. --- 1) 假设 G 是交换群。由 V, p. 390 的推论 7，G 在复希尔伯特空间中的每个不可约酉表示都是一维的。如果 G 局部紧，则集合 \(\widehat{G}\) 可与 G 的酉特征标集合等同（II, p. 201 的定义 1），所以记号 \(\widehat{G}\) 与 II, p. 201 的定义 2 中引入的记号相容。

\begin{enumerate}
\def\labelenumi{\arabic{enumi})}
\setcounter{enumi}{1}
\item
  如果 G 是有限群，则集合 \(\widehat{\mathbf{G}}\) 与 \(\mathbf{C}[\mathrm{G}]\)-单模的类所成的集合之间存在双射（A, VIII, p. 47），后者在 A, VIII, p. 396 中也记作 \(\widehat{\mathbf{G}}\)。
\item
  对于 \(\pi \in \widehat{\mathbf{G}}\)，我们把 \(\overline{\pi}\) 与 \(\widehat{\mathbf{G}}\) 中 \(\pi\) 的共轭表示的类等同。
\item
  如果 \(\pi\) 是 G 在有限维复希尔伯特空间中的不可约酉表示，则其特征标 \(\chi_{\pi}\) 只取决于 \(\pi\) 的类。因此可以谈论 G 的有限维复不可约酉表示的特征标集合。
\end{enumerate}

